\documentclass[t]{beamer}
\usepackage{pgfpages}
\usepackage[ngerman]{babel}
\usepackage{textcomp}
\usepackage{subfigure}
\usepackage{graphicx}
\usepackage{fancyhdr}
\usepackage{tgpagella}
\usepackage{textpos}
\usepackage{hyperref}
\usepackage{tcolorbox}
\usepackage{tikz}
\usetikzlibrary{positioning}

\usepackage{fontspec}
\setmainfont{Titillium Web}
\setsansfont{Titillium Web}

\usetheme{Dresden}
\useoutertheme{infolines}
\usecolortheme{fsfw}

\title{Linux Presentation Day April 2016}
\author{\texttt{https://fsfw-dresden.de/lpd}}
\date[30.04.2016]{30. April 2016}

\newcommand{\cmd}[1]{\texttt{#1}}

\addtobeamertemplate{frametitle}{}{%
    \begin{textblock*}{130mm}(.82\textwidth,-0.5cm)
      \includegraphics[width=1.8cm]{pdf-gen/lpd-logo-small.pdf}
\end{textblock*}
\vspace*{1em}}

\setbeamertemplate{navigation symbols}{}%remove navigation symbols

\begin{document}
\begin{frame}
  \frametitle{Programm}
  \vspace*{-1em}
  \begin{description}
  \item[15:00] Beginn
  \item[15:15] Begrüßung und Vorstellung
  \item[15:30] Einführungsvortrag
  \item[16:15] Pause (Getränke und Gebäck)
  \item[16:30] Praxisbeispiel: Autarke WLAN-Station (IT 4 Refugees)
  \item[17:00] Workshop: Linux-Install-Party
  \item[20:00] Ende
  \end{description}
  \begin{block}{Außerdem}
    \begin{itemize}
    \item Linux ausprobieren (Demonstrationscomputer)
    \item Mate, Wostok und Gebäck
    \item Infomaterial/Stände der beteiligten Initiativen
    \end{itemize}
  \end{block}
\end{frame}

\end{document}
